\documentclass[11pt]{article}
\usepackage[margin=1in]{geometry}
\usepackage{amsmath,amssymb,amsthm}
\usepackage{graphicx}
\usepackage[hidelinks]{hyperref}

\newtheorem{theorem}{Theorem}
\newtheorem{remark}{Remark}

\newcommand{\card}{\operatorname{card}}
\newcommand{\dens}{\operatorname{dens}}
\newcommand{\eps}{\varepsilon}

\begin{document}

\begin{center}
{\Large A separable subcase of the cardinality-injective problem}\\[4pt]
{\normalsize Partial packet for arXiv:1301.3656}
\end{center}

\section*{Source and target}

The source is Lingxin Bao, Lixin Cheng, Qingjin Cheng, and Duanxu Dai,
\emph{On universal left-stability of $\epsilon$-isometries},
arXiv:1301.3656 \cite{BaoChengChengDai}.  In Definition 3.1 the source calls
a Banach space \(X\) \emph{cardinality injective} if there is
\(\lambda<\infty\) such that, whenever \(X\) isometrically embeds into a
Banach space \(Y\) with \(\card(Y)=\card(X)\), the image of \(X\) is the range
of a projection of norm at most \(\lambda\).  The source proves that this is
equivalent to universal left-stability for standard \(\eps\)-isometries and
that dual cardinality-injective spaces are injective.  Remark 3.6 then leaves
the following general question.

\begin{center}
\includegraphics[width=.95\linewidth]{figures/open_problem_crop.png}
\end{center}

\noindent
Thus the target is whether every cardinality-injective Banach space is
injective, in the usual bounded-extension/\(P_\lambda\) sense used by the
source.

\section*{Result}

\begin{theorem}
Let \(X\) be a separable Banach space.  If \(X\) is cardinality injective in
the sense of Bao--Cheng--Cheng--Dai, then \(X\) is finite-dimensional.  Hence,
for separable Banach spaces, the source question has a positive answer:
cardinality injective implies injective.
\end{theorem}

\section*{Proof intuition}

For an infinite-dimensional separable space \(X\), the cardinality restriction
in Definition 3.1 is no restriction on separable superspaces: every
infinite-dimensional separable Banach space has cardinality continuum.
Therefore cardinality injectivity makes \(X\) separably injective.  Zippin's
classification then says that \(X\) must be isomorphic to \(c_0\).  But
\(c_0\) cannot satisfy Definition 3.1, because it is an isometric subspace of
\(\ell_\infty\), \(\card(c_0)=\card(\ell_\infty)\), and Phillips' theorem says
there is no bounded projection \(\ell_\infty\to c_0\).  This rules out the
infinite separable case.  Finite-dimensional spaces are injective by the
standard Hahn--Banach coordinate-extension argument.

\section*{Proof}

Assume first that \(X\) is infinite-dimensional and separable.  Then
\(\card(X)=2^{\aleph_0}\).  If \(Y\) is any infinite-dimensional separable
Banach space containing \(X\), then likewise \(\card(Y)=2^{\aleph_0}\).  Hence
the defining projection property for cardinality injectivity applies to every
separable superspace \(Y\) of \(X\).  Equivalently, \(X\) is separably
injective: every bounded operator from a closed subspace of a separable Banach
space into \(X\) extends to the whole separable space with a uniform bound.

By Zippin's theorem \cite{Zippin}, a separable separably injective Banach
space is isomorphic to \(c_0\).  Proposition 3.2 of
\cite{BaoChengChengDai} says that cardinality injectivity is preserved under
Banach-space isomorphism.  Therefore, if our infinite-dimensional separable
\(X\) were cardinality injective, then \(c_0\) would be cardinality injective.

This is impossible.  The standard inclusion \(c_0\subset \ell_\infty\) is
isometric, and both spaces have cardinality \(2^{\aleph_0}\).  Cardinality
injectivity of \(c_0\) would therefore give a bounded projection
\(\ell_\infty\to c_0\).  Phillips' theorem \cite{Phillips} says that no such
projection exists.  This contradiction proves that no infinite-dimensional
separable cardinality-injective Banach space exists.

It remains only to note that finite-dimensional Banach spaces are injective
in the bounded-extension sense used in the source.  Let \(X\) be
\(n\)-dimensional and isometrically contained in a Banach space \(Y\).  Choose
an Auerbach basis \((x_i,x_i^*)_{i=1}^n\) for \(X\).  Extend each \(x_i^*\) to
\(\phi_i\in Y^*\) by Hahn--Banach with \(\|\phi_i\|=\|x_i^*\|=1\), and set
\[
  P(y)=\sum_{i=1}^n \phi_i(y)x_i \qquad (y\in Y).
\]
Then \(P:Y\to X\) is a projection and \(\|P\|\le n\).  Thus finite-dimensional
spaces are \(P_n\)-spaces, hence injective in the sense of the source.  The
separable case follows.

\section*{Verification and limitations}

The local run indexes were searched for the exact arXiv id, for
\texttt{cardinality injective}, and for universal left-stability keywords.
No existing packet for arXiv:1301.3656 was found.  A neighboring attempt for
arXiv:1301.3396 concerns broader epsilon-isometry stability questions and not
Remark 3.6.  The local later source arXiv:2605.06003 mentions the same
cardinality-injective condition but does not settle this question.  Web
searches for exact phrases around cardinality injectivity found the source
paper and no later exact full resolution during this pass.

This packet proves only the separable subcase.  It does not settle the
nonseparable cardinality-injective versus injective problem.

\section*{Human-review recommendation}

Likely valid partial result.  Review should check that the source's
``injective'' convention is the bounded \(P_\lambda\)/extension convention,
that Proposition 3.2 is being used with the same isomorphic notion, and that
the cited Zippin and Phillips theorems are acceptable standard inputs.

\begin{thebibliography}{9}

\bibitem{BaoChengChengDai}
L. Bao, L. Cheng, Q. Cheng, and D. Dai,
\emph{On universal left-stability of \(\epsilon\)-isometries},
arXiv:1301.3656; Acta Math. Sin. (Engl. Ser.) 29 (2013), no. 11,
2037--2046.

\bibitem{Phillips}
R. S. Phillips,
\emph{On linear transformations},
Trans. Amer. Math. Soc. 48 (1940), 516--541.

\bibitem{Zippin}
M. Zippin,
\emph{The separable extension problem},
Israel J. Math. 26 (1977), 372--387.

\end{thebibliography}

\end{document}
